% Lösungen Einheit 2 – Prozentwert berechnen
\einheitenkopf*{Einheit 2 von 4 · Prozentwert berechnen}
\erg{22}{Beispiel 15\,€ \quad a) 60\,kg \quad b) 30\,kg \quad c) 15\,kg \quad d) 6\,kg \quad e) Das Ganze (60\,kg) bleibt gleich; die Prozentzahl wird kleiner, dann wird auch der Teil kleiner.}
\erg{23}{Beispiel 6\,€ \quad a) 1\,\% $=$ 3\,€, 7\,\% $=$ 21\,€ \quad b) 10\,\% $=$ 9\,€, 30\,\% $=$ 27\,€ \quad c) 10\,\% $=$ 6\,€, 5\,\% $=$ 3\,€, 15\,\% $=$ 9\,€}
\erg{24}{Beispiel 14\,€ \quad a) $0{,}6 \cdot 45 = 27$\,kg \quad b) $0{,}3 \cdot 2{,}50 = 0{,}75$\,€}
\erg{25}{a) 15\,€ \quad b) $60 - 15 = 45$\,€ \quad c) $1{,}5 \cdot 60 = 90$ Mitglieder}
\erg{26}{a) $0{,}19 \cdot 300 = 57$\,€ \quad b) 54\,€ ($0{,}15 \cdot 360$; 306\,€ ist der neue Preis)}
\erg{27}{a) Pia schreibt 40\,\% als 0,04 statt 0,4; richtig $0{,}4 \cdot 25 = 10$\,€ \quad b) $0{,}8 \cdot 35 = 28$\,€}
\erg{28}{a) 63\,€ ist der neue Preis, nicht die Ersparnis; man spart 7\,€ \quad b) $0{,}15 \cdot 80 = 12$\,€}
\erg{29}{a) 1\,\% ist der hundertste Teil des Ganzen: $400 : 100 = 4$\,€; dann mal 9: 36\,€. \quad b) Ja: 5\,\% von 200\,€ sind 10\,€, 10\,\% von 90\,€ sind nur 9\,€ – es kommt auch auf das Ganze an.}
\erg{30}{a) 20\,€ \quad b) Nord: $80 - 20 = 60$\,€, Süd: $80 - 15 = 65$\,€ \quad c) Laden Nord ist um 5\,€ günstiger.}
